% ============================================================
%  RoboMME RL Experiment Report
%  CS224R Final Project — Stanford, Spring 2026
%  Compiled 2026-05-28
% ============================================================
\documentclass[11pt,letterpaper]{article}

% --- Packages -----------------------------------------------
\usepackage[margin=1in]{geometry}
\usepackage{amsmath,amssymb}
\usepackage{booktabs}
\usepackage{tabularx}
\usepackage{multirow}
\usepackage{xcolor}
\usepackage{hyperref}
\usepackage{enumitem}
\usepackage{graphicx}
\usepackage{float}
\usepackage{caption}
\usepackage{subcaption}
\usepackage{array}
\usepackage{listings}
\usepackage{microtype}
\usepackage{parskip}
\usepackage{titlesec}
\usepackage{fancyhdr}
\usepackage{mdframed}

% --- Colors -------------------------------------------------
\definecolor{darkblue}{RGB}{10,50,120}
\definecolor{darkred}{RGB}{150,20,20}
\definecolor{darkgreen}{RGB}{20,100,40}
\definecolor{lightgray}{RGB}{245,245,245}
\definecolor{warningbg}{RGB}{255,248,220}

% --- Hyperref -----------------------------------------------
\hypersetup{
  colorlinks=true,
  linkcolor=darkblue,
  citecolor=darkgreen,
  urlcolor=darkblue,
  pdftitle={RoboMME RL Experiments — CS224R},
  pdfauthor={Jevon Mao et al.}
}

% --- Header / Footer ----------------------------------------
\pagestyle{fancy}
\fancyhf{}
\rhead{\textcolor{gray}{\small RoboMME RL Experiments}}
\lhead{\textcolor{gray}{\small CS224R — Spring 2026}}
\rfoot{\thepage}

% --- Section formatting -------------------------------------
\titleformat{\section}{\large\bfseries\color{darkblue}}{}{0em}{}[\titlerule]
\titleformat{\subsection}{\normalsize\bfseries\color{darkblue}}{}{0em}{}
\titleformat{\subsubsection}{\normalsize\bfseries}{}{0em}{}

% --- Code listings ------------------------------------------
\lstset{
  basicstyle=\ttfamily\footnotesize,
  backgroundcolor=\color{lightgray},
  frame=single,
  breaklines=true,
  columns=fullflexible,
  keepspaces=true
}

% --- Custom environments ------------------------------------
\newmdenv[
  backgroundcolor=warningbg,
  linecolor=orange!60!black,
  linewidth=1pt,
  leftmargin=0pt,
  rightmargin=0pt,
  innerleftmargin=8pt,
  innerrightmargin=8pt,
  innertopmargin=6pt,
  innerbottommargin=6pt,
  skipabove=6pt,
  skipbelow=6pt
]{finding}

\newmdenv[
  backgroundcolor=lightgray,
  linecolor=gray!60,
  linewidth=0.5pt,
  leftmargin=0pt,
  rightmargin=0pt,
  innerleftmargin=8pt,
  innerrightmargin=8pt,
  innertopmargin=6pt,
  innerbottommargin=6pt,
  skipabove=6pt,
  skipbelow=6pt
]{codeblock}

% ============================================================
\begin{document}

% --- Title page ---------------------------------------------
\begin{titlepage}
  \centering
  \vspace*{1.5cm}
  {\LARGE\bfseries\color{darkblue} RoboMME Reinforcement Learning Experiments}\\[0.4em]
  {\large\color{gray} BinFill and PickXtimes — Full Experimental Trace}\\[2em]
  {\normalsize CS224R Final Project $\cdot$ Stanford University $\cdot$ Spring 2026}\\[1em]
  {\normalsize Compiled: 2026-05-28}\\[3em]

  \begin{mdframed}[backgroundcolor=lightgray,linecolor=gray!50]
  \textbf{Document purpose.}
  This document is written for a project teammate who is joining at
  the midpoint and needs a complete picture of what has been tried,
  why each approach was chosen, what failed and why, and where the
  strongest next steps lie. It covers every training run from the
  initial BinFill PPO baseline through the seven-phase PickXtimes
  ablation (BinFill: §\ref{sec:binfill}; PickXtimes: §\ref{sec:pickxtimes}).
  All numbers are deterministic-eval success rates on held-out
  test-split episodes unless stated otherwise.
  \end{mdframed}

  \vfill
  {\small\textit{Generated from \texttt{REPORT.md}, \texttt{PICKXTIMES\_REPORT.md},
  \texttt{EXPERIMENT\_LOG.md}, and \texttt{NEXT\_STEPS.md}.}}
\end{titlepage}

\tableofcontents
\newpage

% ============================================================
\section{Background and Task Descriptions}
\label{sec:background}

\subsection{RoboMME Benchmark}

RoboMME is a manipulation benchmark testing memory and multi-step reasoning
in robotic agents. A 7-DoF Panda arm with a parallel-jaw gripper operates in
a physics simulation (SAPIEN / ManiSkill). The official evaluation criterion
is a binary success rate averaged over multiple seeds and tasks.

The benchmark ships with imitation-learning (IL) baselines only — no RL
numbers exist in the paper. Reference IL numbers for context:

\begin{table}[H]
\centering
\begin{tabular}{lc}
\toprule
\textbf{Method} & \textbf{BinFill SR} \\
\midrule
Human & 96.0\% \\
GroundSG + Oracle (ground-truth info) & 85.8\% \\
FrameSamp + Modul (best non-oracle) & 39.6\% \\
\bottomrule
\end{tabular}
\caption{Published RoboMME paper baselines (IL, not RL).}
\end{table}

The closest published RL comparison is \textbf{MIKASA-Robo} (ICLR 2025), which
benchmarks PPO-MLP, PPO-LSTM, SAC, and TD-MPC2 on 32 ManiSkill memory-manipulation
tasks. Their key finding: \emph{``none of the models were able to successfully
solve the majority of MIKASA-Robo tasks.''} PPO and PPO-LSTM collapsed to $\approx$0\%
on harder memory tasks at standard training budgets (up to $\sim$5M steps).

\subsection{Task: BinFill}
\label{sec:binfill-task}

A cluttered-scene manipulation task. The agent must:
\begin{enumerate}[nosep]
  \item Pick up multiple target cubes in a cluttered bin,
  \item Place each in the target bin, and
  \item Press a button to signal completion.
\end{enumerate}
BinFill is among the harder RoboMME tasks due to scene clutter.
Episode horizon: \textbf{1500 steps} (official). Training initially used 500 steps
— a significant mismatch identified post-training.

\subsection{Task: PickXtimes}
\label{sec:pickxtimes-task}

A multi-step pick-and-place task with a counting requirement:
\begin{enumerate}[nosep]
  \item Pick up a specified target cube \textbf{N times} ($N = 1$–$5$, sampled
        per episode by \texttt{seed \% 3}: easy / medium / hard difficulty).
  \item After each pickup, place it on a target pad.
  \item Press a button to stop.
\end{enumerate}

Total subgoals per episode: $2N + 1$ (range: 3–11).
Episode horizon: \textbf{1500 steps}.
Failure conditions: picking the wrong cube or pressing the button before completing
all cycles.

\paragraph{Observation space.}
128$\times$128 front + wrist RGB images $\oplus$ 7-D joint state $\oplus$ 6-D EEF
state $\oplus$ 2-D gripper state (\texttt{obs\_mode=rgb}).

\paragraph{Action space.}
8-D continuous in $[-1, 1]$: 7 joint $q_\text{pos}$ targets $+$ 1 gripper command.
\texttt{normalize\_action=False} — $[-1, 1]$ is the literal joint-target range,
clamped by the wrapper before stepping ManiSkill.

\paragraph{Why PickXtimes?}
It isolates \emph{both} memory and curiosity axes cleanly:
\begin{itemize}[nosep]
  \item \textbf{Memory}: the agent must count pickups — a purely reactive policy
        caps at $1/N$ success at best.
  \item \textbf{Curiosity}: sparse extrinsic reward requires the agent to discover
        the repeated pick primitive.
  \item \textbf{Motor}: a single object type and one primitive — moderate difficulty,
        not the primary bottleneck.
\end{itemize}

% ============================================================
\section{Shared Infrastructure}

\subsection{Hardware and Software}

\begin{table}[H]
\centering
\begin{tabular}{ll}
\toprule
\textbf{Component} & \textbf{Details} \\
\midrule
Cluster & Stanford SC — partitions \texttt{svl}/\texttt{svl-interactive}/\texttt{visionlab} \\
GPU & NVIDIA Titan RTX (24 GB) \\
Local dev & NVIDIA RTX 4090, WSL2 (code dev only; SAPIEN smoke tests on cluster) \\
Python & 3.11 \\
PyTorch & 2.9.1+cu121 \\
SB3 / sb3-contrib & stable\_baselines3 2.8.0 / sb3-contrib 2.8.0 \\
ManiSkill & 3.0.0b21 (YinpeiDai fork, git rev \texttt{07be6fbc}) \\
SAPIEN & 3.0.2 \\
\midrule
\multicolumn{2}{l}{\textbf{Required env var:} \texttt{ROBOMME\_SIM\_BACKEND=physx\_cpu}} \\
\multicolumn{2}{l}{SAPIEN 3.0.3 + torch 2.9 trip a \texttt{\_\_cuda\_array\_interface\_\_} typestr} \\
\multicolumn{2}{l}{parse error in \texttt{articulation\_joint.drive\_target} on GPU sim.} \\
\multicolumn{2}{l}{CPU sim runs at 60–75 FPS aggregate at \texttt{n\_envs=4} on Titan RTX.} \\
\bottomrule
\end{tabular}
\caption{Environment. All cluster jobs export \texttt{ROBOMME\_SIM\_BACKEND=physx\_cpu}.}
\end{table}

\subsection{Shared Visual Backbone}

All policies share a \textbf{ResNet-18 CNN feature extractor} (\texttt{RobommeCNNExtractor})
with per-camera projection heads outputting a 576-dimensional joint feature vector,
shared between actor and critic.

\subsection{Policy Architecture Variants}

\begin{description}[leftmargin=2em]
  \item[PPO (vanilla)] MLP head on the 576-d extractor features.
  \item[PPO+ICM] Same as PPO + Pathak ICM: forward-dynamics MSE as intrinsic reward.
        Intrinsic return normalized by running std (Burda 2018 recipe).
  \item[RecurrentPPO (LSTM)] 576-d features $\to$ LSTM (hidden=256, 1 layer, shared
        actor/critic) $\to$ policy/value MLPs. Implemented via \texttt{sb3-contrib}
        \texttt{MultiInputLstmPolicy}.
  \item[RND] Predictor + frozen-random target, both 4-layer MLP heads on 576-d
        features. Intrinsic reward $= \|\text{predictor}(\phi) - \text{target}(\phi)\|^2$,
        normalized by running std of intrinsic returns. Predictor gradient does
        \emph{not} flow into the shared extractor.
  \item[RecurrentPPO+RND] Combines LSTM memory with RND intrinsic reward.
\end{description}

% ============================================================
\section{BinFill Experiments}
\label{sec:binfill}

\subsection{Training Configurations}

\begin{table}[H]
\centering
\begin{tabular}{lll}
\toprule
\textbf{Run} & \textbf{Baseline} & \textbf{Key config} \\
\midrule
\texttt{ppo\_v2} & Vanilla PPO (original reward stub) & 500 steps, reward stub \\
\texttt{ppo\_v3} & Vanilla PPO (v3 shaping) & 500 steps, negative shaping \\
\texttt{ppo\_v4} & Vanilla PPO (v4 shaping) & 500 steps, positive proximity shaping \\
\texttt{ppo\_icm\_v4} & PPO+ICM (v4 reward) & 500 steps, killed at 88k \\
\texttt{ppo\_icm\_v7} & PPO+ICM (v4 reward + opts) & sim\_freq=40, camera 128×128 \\
\bottomrule
\end{tabular}
\caption{BinFill training runs. All use \texttt{n\_envs=4}, \texttt{n\_steps=512},
\texttt{batch\_size=256}, \texttt{lr=3e-4}.}
\end{table}

\paragraph{Shaped reward (v4).}
\[
r_t = 0.05 \cdot \frac{1}{1 + d_t} + r_\text{task}
\]
where $d_t$ is TCP-to-cube distance and $r_\text{task}$ includes per-pickup bonuses
and a terminal success bonus.

\subsection{Training Progression}

\begin{table}[H]
\centering
\begin{tabular}{llllll}
\toprule
\textbf{Run} & \textbf{Steps} & \textbf{Final \texttt{ep\_rew\_mean}} &
\textbf{Final \texttt{ep\_len\_mean}} & \textbf{Wall-clock} & \textbf{FPS} \\
\midrule
\texttt{ppo\_v2}     & 500k & 0.0   & 501 & 4.8 h & 28 \\
\texttt{ppo\_v3}     & 252k & $-5.1$ & 501 & 2.4 h & 29 \\
\texttt{ppo\_v4}     & 252k & 10.3  & 501 & 2.4 h & 29 \\
\texttt{ppo\_icm\_v4} & 88k (killed) & 10.4 & 501 & 1.5 h & 16 \\
\texttt{ppo\_icm\_v7} & 254k & 11.5 & 501 & 1.0 h & 75 \\
\bottomrule
\end{tabular}
\caption{BinFill training summary. \texttt{ep\_len\_mean = 501} across every run
indicates no episode ever terminated via success during training; the 500-step
horizon was always exhausted.}
\end{table}

\begin{finding}
\textbf{Key observation.} \texttt{ep\_len\_mean = 501} throughout all BinFill runs.
The policy is climbing the shaped reward (approach cubes, partial grasps) but
never closes out a full BinFill episode within the 500-step training horizon.
The +12\% gap ICM showed during training (10.3 $\to$ 11.5) did not translate
into any task completion.
\end{finding}

\subsection{Evaluation Results}

After identifying the 500-step training horizon mismatch, evaluations were run
at both 500 and 1500 steps (official).

\begin{table}[H]
\centering
\begin{tabular}{llccc}
\toprule
\textbf{Run} & \textbf{Horizon} & \textbf{Episodes} & \textbf{Success rate} &
\textbf{Notes} \\
\midrule
\texttt{ppo\_v4}     & 500  & 20 & 0.0\% & training-matched \\
\texttt{ppo\_icm\_v7} & 500 & 20 & 0.0\% & training-matched \\
\texttt{ppo\_v4}     & 1500 & 20 & \textbf{0.0\%} & official; sim\_freq=100, cam=256 \\
\texttt{ppo\_icm\_v7} & 1500 & 20 & \textbf{0.0\%} & physics shift vs training \\
\bottomrule
\end{tabular}
\caption{BinFill deterministic eval results, single seed (seed=0).}
\end{table}

\subsection{Configuration Caveats}
\label{sec:binfill-caveats}

\begin{enumerate}
  \item \textbf{Horizon mismatch.} Training at 500 steps; official eval at 1500.
        Multi-stage tasks may be impossible in 500 steps regardless of policy quality.
        Fixed in code; not yet retrained.
  \item \textbf{Physics mismatch (\texttt{ppo\_icm\_v7}).} Trained at
        \texttt{sim\_freq=40} (40 Hz, 2 substeps/env.step); official is
        \texttt{sim\_freq=100} (5 substeps). Trained-at-coarse, eval-at-coarse is
        internally consistent but introduces distribution shift vs official eval.
        Default removed for future runs.
  \item \textbf{Camera mismatch (\texttt{ppo\_icm\_v7}).} Trained with cameras at
        128$\times$128; official renders at 256. Minor distribution shift. Removed
        for future runs.
  \item \textbf{Single seed.} Seed=0 only; official requires seeds \{7, 42, 0\} averaged.
  \item \textbf{Single task.} BinFill is one of 16 benchmark tasks and one of the
        hardest (cluttered scenes).
\end{enumerate}

% ============================================================
\section{PickXtimes Experiments}
\label{sec:pickxtimes}

\subsection{Common Training Configuration}

\begin{table}[H]
\centering
\begin{tabular}{lll}
\toprule
\textbf{Param} & \textbf{Value} & \textbf{Notes} \\
\midrule
\texttt{n\_envs} & 4 & SubprocVecEnv \\
\texttt{n\_steps} & 512 & rollout horizon per env \\
\texttt{batch\_size} & 256 & PPO minibatch \\
\texttt{n\_epochs} & 4 & PPO update epochs per rollout \\
\texttt{lr} & 3e-4 & policy/value optimizer \\
\texttt{gamma} & 0.997 & high discount for sparse reward \\
\texttt{target\_kl} & 0.05 & early-stop PPO update \\
\texttt{ent\_coef} & 0.01 & (varied to 0.001 in one sweep) \\
\texttt{gae\_lambda} & 0.95 & default \\
\texttt{max\_steps} & 1500 & matches \texttt{phase1\_eval.py} \\
\texttt{LSTM hidden} & 256 & 1 shared layer (RecurrentPPO runs) \\
\texttt{RND output dim} & 128 & 4-layer MLP predictor + target \\
\bottomrule
\end{tabular}
\caption{Common PickXtimes training hyperparameters.}
\end{table}

\subsection{Reward Variants}
\label{sec:reward-variants}

All variants share a step penalty $r_\text{step} = -0.01$.

\begin{table}[H]
\centering
\renewcommand{\arraystretch}{1.2}
\begin{tabular}{lccccccc}
\toprule
& \textbf{Prox.\ coef} & \textbf{Subgoal bonus} & \textbf{Pickup bonus}
& \textbf{Grasp hint} & \textbf{Succ.\ terminal} & \textbf{Fail terminal} \\
\midrule
\textbf{V1} & 0.05 & +2.0 / subgoal & +1.0 / pick & +0.05 / step & +10 & $-5$ \\
\textbf{V2} & 0.005 & +10 & +5 & +0.5 & +50 & $-10$ \\
\textbf{V3} & 0.05 (V1) & +10 (V2) & +5 (V2) & +0.5 (V2) & +50 & $-10$ \\
\textbf{V4} & \textbf{0} & +10 & +5 & +0.5 & \textbf{+100} & $-10$ \\
\bottomrule
\end{tabular}
\caption{PickXtimes reward variants. V1 is proximity-exploitable; V4 removes
proximity entirely to force genuine task progress.}
\end{table}

\paragraph{Proximity reward formula.}
\[
r_\text{prox} = \frac{\alpha}{1 + d_t}
\]
where $d_t$ is the TCP-to-target-cube distance and $\alpha$ is the
\emph{Prox.\ coef} from Table~4. In V4, $\alpha = 0$.

\subsection{Experimental Matrix}

\begin{table}[H]
\centering
\renewcommand{\arraystretch}{1.15}
\begin{tabular}{lcccccl}
\toprule
\textbf{Run label} & \textbf{Memory} & \textbf{Curiosity} & \textbf{BC warm-start}
& \textbf{Reward} & \textbf{Seeds} & \textbf{Final SR} \\
\midrule
b1 V1 (PPO) & \texttimes & \texttimes & \texttimes & V1 & 1 killed & 0/10 @ 900k \\
b4 V1 (PPO+RND) & \texttimes & \checkmark & \texttimes & V1 & 1 killed & inferred 0\% \\
b5 V1 (RecurrentPPO) & \checkmark & \texttimes & \texttimes & V1 & 1 $\to$ 1.34M & 0/20 @ 700k \\
b6 V1 (Recurrent+RND) & \checkmark & \checkmark & \texttimes & V1 & 3 $\to$ 1.17M & 0/20 @ 900k \\
b6 V1 + BC & \checkmark & \checkmark & \checkmark & V1 & 3 $\sim$1M & 0/20 @ 600k \\
b1–b6 V2 (all cells) & all & & \texttimes & V2 & 3 @ $\sim$700k & 0/20 @ 600k \\
b6 V3 + BC & \checkmark & \checkmark & \checkmark & V3 & 3 @ $\sim$560k & 0/20 @ 500k \\
b6 V3 + BC + ent=0.001 & \checkmark & \checkmark & \checkmark & V3+low ent & 3 @ $\sim$640k & 0/30 subgoals \\
V4 + BC + aux=0.5 & \checkmark & \checkmark & \checkmark & V4 & 3 $\to$ 5M & stuck at $-15$ floor \\
DAgger v1 (4 iters) & (BC retrain) & — & \checkmark & — & 1 & 0/20 pickups \\
DAgger v2 (8 iters) & (BC retrain) & — & \checkmark & — & 1 & 0/20 through iter 5 \\
DAPG (bc\_aux=0.05) & \checkmark & \checkmark & \checkmark & V4 & 3 killed @590k & stuck at $-15$ \\
\bottomrule
\end{tabular}
\caption{Full PickXtimes experimental matrix. SR = test-split success rate.
``Subgoals'' = max subgoal index reached (0 means first pickup never completed).}
\end{table}

\noindent
\textbf{Total cluster compute:} ${\approx}33$ RL jobs $\times$ ${\approx}5$h
$\approx$ \textbf{165 Titan RTX-hours.}

% ============================================================
\subsection{Phase 1: Initial Reward Sweep (V1, $\sim$1.5M steps)}

\paragraph{Configuration.}
12 jobs: $\{$b1 PPO, b4 PPO+RND, b5 RecurrentPPO, b6 RecurrentPPO+RND$\}
\times$ 3 seeds $\times$ V1 reward $\times$ 1.5M steps.

\paragraph{Results.}

\begin{table}[H]
\centering
\begin{tabular}{lccccc}
\toprule
\textbf{Variant} & \textbf{Best train rew.} & \textbf{SR @ 300k}
& \textbf{SR @ 500k} & \textbf{SR @ 700k} & \textbf{SR @ 900k} \\
\midrule
b1 (PPO) & 35.9 & 0/10 & 0/10 & — & 0/10 \\
b4 (PPO+RND) & 31.4 & 0/10 & — & — & — \\
b5 (RecurrentPPO) & 46.6 & 0/10 & 0/10 & 0/20 & 0/20 \\
b6 (RecurrentPPO+RND) & 40.4 & 0/10 & — & — & 0/20 \\
\bottomrule
\end{tabular}
\caption{Phase 1 results. b5 subgoal diagnostic at 900k: 0/20 episodes reach
even subgoal 1 (first pickup).}
\end{table}

Training reward climbed smoothly 25 $\to$ 46 over $\sim$900k steps \emph{without
spikes} — characteristic of proximity-hacking (``hover near object'') rather
than subgoal completion.

\subsection{Phase 2: Behavior Cloning Pipeline}

\paragraph{Data collection.}
For each of 80 train-split episodes, \texttt{PandaArmMotionPlanningSolver} was
instantiated against the env and each subgoal's \texttt{solve(env, planner)}
callback was executed. \texttt{env.step} was hooked to record
\texttt{(wrapped\_obs, action)} pairs as per-episode \texttt{.npz} files.

\begin{table}[H]
\centering
\begin{tabular}{lll}
\toprule
\textbf{Dataset} & \textbf{Episodes / Transitions} & \textbf{BC loss} \\
\midrule
Expert (non-recurrent) & 80 ep, 42{,}821 transitions & 0.006 \\
Expert (recurrent) & 80 ep, 42{,}821 transitions & 0.007 \\
\bottomrule
\end{tabular}
\end{table}

\paragraph{BC training.}
Same policy class as PPO; MSE between \texttt{policy\_mean} and expert action.
For the recurrent variant: per-sample LSTM init-zero state (treat each
transition as memoryless). 30 epochs, batch 256, lr 3e-4.

\paragraph{BC-only evaluation.}
\begin{itemize}[nosep]
  \item Test split, deterministic: \textbf{0/20}
  \item Train split (same episodes used for BC training), deterministic:
        \textbf{0/20}
\end{itemize}

Even on the same episodes the BC model was trained on, the BC policy never
completes the first pickup. Compounding error over 1500 steps puts the policy
off the expert manifold within $\sim$50 steps.

\subsection{Phase 3: BC Warm-Start Sweep}

The BC-trained \texttt{state\_dict} was loaded into the PPO policy at startup,
then RL training continued.

\begin{table}[H]
\centering
\begin{tabular}{lll}
\toprule
\textbf{Config} & \textbf{Reward} & \textbf{SR @ 600k} \\
\midrule
b6 + BC & V1 & 0/20 \\
b6 + BC & V2 & train reward $\approx -9.6$ (barely above $-10$ floor) \\
b6 + BC & V3 & 0/20 \\
b6 + BC + \texttt{ent\_coef=0.001} & V3 & 0/30 (subgoal diagnostic, 3 seeds) \\
\bottomrule
\end{tabular}
\end{table}

BC warm-start yielded a +2 training-reward bump (b6+BC $\approx$ 43 vs b6 $\approx$ 41)
but no improvement in eval SR. Reducing \texttt{ent\_coef} from 0.01 to 0.001
lowered policy std from 1.02 to 0.83 but produced the same 0/30 subgoal result.

\subsection{Phase 4: V4 Reward + BC + \texttt{bc\_aux\_coef=0.5} (5M steps)}

\paragraph{Design intent.}
Remove proximity reward entirely ($\alpha = 0$); rely on auxiliary BC loss to
anchor the policy near the expert mean during RL training.

\paragraph{Configuration.}
3 seeds $\times$ 5M step target; \texttt{bc\_aux\_coef=0.5} (auxiliary
\texttt{MSE(policy\_mean, expert\_action)} loss applied after each PPO update
using a 30k-transition demo subset).

\paragraph{Result.}
Stuck at training-reward floor of $-15$ (V4's pure step penalty $\times$ 1500)
for 850k steps before killed. Policy std dropped from 0.37 $\to$ 0.55.
Zero subgoal completions across all 3 seeds.

\begin{finding}
\textbf{Root cause.} \texttt{bc\_aux\_coef=0.5} was \textbf{10$\times$ too strong}.
The auxiliary BC loss pinned the policy to the BC mean (which itself was 0\%
capable on its own). PPO's policy gradient was effectively suppressed; the run
degraded to ``what BC alone achieves, maintained by RL.''
\end{finding}

\subsection{Phase 5: DAgger v1 (4 Iterations, $K=15$)}

\paragraph{Expert-query mechanism.}
At every $K=15$ env-steps during rollout, the current environment state is
passed to the motion planner (\texttt{PandaArmMotionPlanningSolver}). The first
action the planner would take from that state is captured as the expert label.

\begin{codeblock}
\begin{lstlisting}[language=python]
def get_first_expert_action(env, planner):
    # Patch open/close gripper to no-ops (6-step stay-put prelude)
    planner.open_gripper  = noop_open
    planner.close_gripper = noop_close
    def hooked(action, *a, **kw):
        captured.append(np.clip(action, -1, 1))
        raise _StopAfterFirstStep()
    env.step = hooked
    try:
        solve_fn(env, planner)   # motion planner runs here
    except (_StopAfterFirstStep, Exception):
        pass
    return captured[0] if captured else None
\end{lstlisting}
\end{codeblock}

\paragraph{Bug (iters 1–2).}
The hook was capturing the first \texttt{env.step} call, but
\texttt{solve\_pickup} starts with \texttt{planner.open\_gripper(t=6)}, emitting
6 ``stay-put + gripper-open'' no-op actions before any motion. Iters 1–2
collected $\sim$4000 no-op transitions (100\% query success on these trivial
actions), polluting the dataset.

\paragraph{Fix (iters 3–4).}
Patch \texttt{planner.open\_gripper}/\texttt{close\_gripper} to no-ops before
invoking \texttt{solve}. Query success dropped to 52–67/100, indicating real
motion planning.

\begin{table}[H]
\centering
\begin{tabular}{lccc}
\toprule
\textbf{Iter} & \textbf{Aggregated transitions} & \textbf{Clean fraction} & \textbf{Inline eval} \\
\midrule
1 (BC baseline) & 44,821 & — & (eval bug) \\
2 (polluted) & 46,821 & 0\% extra & (eval bug) \\
3 (clean) & 48,031 & 4.2\% extra & (eval bug) \\
4 (clean) & 48,551 & 4.4\% extra & \textbf{0/10 pickups} \\
Final eval & — & — & \textbf{0/20 pickups, 0/20 drops} \\
\bottomrule
\end{tabular}
\end{table}

\begin{finding}
\textbf{Failure mode.} Polluted transitions (40/160 = 25\% of episodes) swamped the
clean DAgger corrections (only 2,000 / 48,551 = 4\% of final dataset). The dataset
was too polluted to show improvement.
\end{finding}

\subsection{Phase 6: DAgger v2 (8 Iterations, $K=5$, Fresh Dataset)}

Fresh aggregated directory (no polluted-data carryover). $K=5$ (3$\times$ more
expert queries per episode), 8 iters, targeting $\sim$24,000 clean DAgger
transitions on top of 42,821 BC transitions ($\approx$36\% clean fraction at
completion).

\begin{table}[H]
\centering
\begin{tabular}{lccc}
\toprule
\textbf{Iter} & \textbf{Aggregated transitions} & \textbf{BC retrain loss} & \textbf{Policy SR} \\
\midrule
1 & 47,282 & $\sim$0.007 & 0/20 subgoals \\
2 & 50,667 & $\sim$0.007 & 0/20 subgoals \\
3 & 53,821 & $\sim$0.007 & 0/20 subgoals \\
4 & 56,618 & $\sim$0.007 & 0/20 subgoals \\
5 & 60,308 & $\sim$0.007 & 0/20 subgoals \\
6–7 & OOM'd (pickle 4-GiB limit) & — & — \\
\bottomrule
\end{tabular}
\end{table}

\begin{finding}
\textbf{Failure mode.} BC retraining consistently converged to loss $\approx$0.007
regardless of DAgger iteration. The per-state MSE objective appears fundamentally
unable to capture trajectory-level success at our model capacity and training
scale. Even with 25\% clean correction data, the resulting policy cannot complete
the first pickup. The job was also killed at iter 7 due to a pickle protocol 3
$\to$ 4 GiB string limit OOM.
\end{finding}

\subsection{Phase 7: DAPG-Style Enhanced PPO (5M Steps $\times$ 3 Seeds)}

\paragraph{Design.}
This run attempted to correct every identified failure mode simultaneously:

\begin{table}[H]
\centering
\begin{tabular}{ll}
\toprule
\textbf{Element} & \textbf{Setting} \\
\midrule
Reward & V4 (no proximity hack) \\
Memory & RecurrentPPO (LSTM 256) \\
Curiosity & RND (128-d, running RMS normalization) \\
BC warm-start & DAgger v2 iter-5 policy \\
\texttt{init\_log\_std} & $-1.0$ (std = 0.37, tight exploration) \\
\texttt{bc\_aux\_coef} & \textbf{0.05} (10$\times$ lower than Phase 4) \\
BC aux dataset & 30k transitions \\
Step budget & 5M $\times$ 3 seeds (chained A+B jobs) \\
\bottomrule
\end{tabular}
\end{table}

\paragraph{Result at 590k steps (12\% through).}

\begin{table}[H]
\centering
\begin{tabular}{lccccc}
\toprule
\textbf{Seed} & \textbf{Steps} & \textbf{\texttt{ep\_rew\_mean}} & \textbf{std}
& \textbf{approx\_kl} & \textbf{value\_loss} \\
\midrule
7  & 589k & $-15.0$ (floor) & 0.513 & 0.025 & 0.02 \\
42 & 588k & $-15.0$ (floor) & 0.512 & 0.025 & 0.02 \\
0  & 588k & $-15.0$ (floor) & 0.514 & 0.025 & 0.02 \\
\bottomrule
\end{tabular}
\end{table}

\texttt{value\_loss} $\approx$ 0.02 means the value function trivially fits
the constant $-15$ reward. \texttt{approx\_kl} well below \texttt{target\_kl=0.05}
— PPO doesn't even early-stop. Std rising 0.37 $\to$ 0.51 — PPO widening
exploration because it has no signal to commit to.

\begin{finding}
\textbf{Decision: terminated} at 590k steps.
Probability of breakthrough at full 5M judged $<$10\% given no signal across
all 3 seeds at the same point where the Phase~4 run had already shown the same
flat-floor pattern. Continuing would waste $\sim$75 GPU-hours.
\end{finding}

% ============================================================
\section{Diagnosis of Failures}
\label{sec:diagnosis}

The 7-phase ablation isolates the failure to two co-occurring structural problems
that cannot be fixed by reward / memory / curiosity / BC tweaks at the scales tested:

\subsection{F1 — Proximity Shaping is Reward-Hackable}

V1's $0.05/(1+d)$ per-step term pays up to $\sim$50–70 over a 1500-step episode
for pure ``hover near object'' behavior, dwarfing the +14 subgoal + +10 terminal
bonuses. Training reward climbed smoothly without spikes. Diagnostic: 0/20
episodes reach subgoal 1 at 900k steps despite training reward of 46.

\subsection{F2 — Removing Proximity Removes the Only Exploration Gradient}

V2 attenuated proximity 10$\times$ and boosted task bonuses 5$\times$. With no
gradient to guide the policy toward cubes, random exploration in 8-D continuous
action space (std $\approx$ 1) never produces a coordinated grasp sequence.
Training reward stayed at floor ($-10$) for the full 1.1M steps tested.

\subsection{F3 — BC Compounding Errors over 1500 Steps}

BC loss converges to 0.006 (excellent per-state fitting) but the resulting
policy fails 100\% even on TRAIN-split episodes. Per-state MSE cannot prevent
trajectory-level compounding errors over 1500 steps; the policy drifts off the
expert manifold within $\sim$50 steps and never recovers.

\subsection{F4 — RND Curiosity Solves the Wrong Exploration Problem}

RND rewards visiting novel CNN-extractor feature states. Our exploration
challenge is ``execute precise gripper close at the right height.'' RND has no
signal for fine motor precision. RND on top of RecurrentPPO underperformed
RecurrentPPO alone in every run (b6: 40 vs b5: 46 at 600k).

\subsection{F5 — LSTM Memory Helps Training Reward but Not Eval}

RecurrentPPO reaches training reward 46 vs vanilla PPO 36. LSTM clearly aids
tracking the target across the episode. However, memory has no job on subgoal 0
(which the policy never passes); counting picks never matters.

\subsection{F6 — High Exploration Noise Drowns Out BC Warm-Start}

With default \texttt{init\_log\_std=0} (std=1.0), the BC-loaded policy's near-expert
mean is immediately dominated by Gaussian noise. The first sampled action deviates
from BC mean by $\sim\sigma$ per dimension — far too much for the fine-grained
grasp configuration.

\subsection{F7 — Auxiliary BC Loss Magnitude is Fragile}

\texttt{bc\_aux\_coef=0.5} (Phase 4) pinned the policy to BC mean (which itself was
0\% capable), suppressing PPO's gradient. The run stayed at $-15$ floor for 850k
steps. \texttt{bc\_aux\_coef=0.05} (Phase 7) is 10$\times$ lower but still
insufficient — PPO still has no positive signal to commit to in V4.

\subsection{F8 — DAgger Bottlenecked by BC Retraining Objective}

Both DAgger v1 (polluted) and v2 (clean) produce 0\% policy success at each
iteration. BC loss converges to 0.007 — virtually identical to plain BC.
The per-state MSE objective fundamentally cannot capture trajectory-level success
at our policy class and training scale.

% ============================================================
\section{Summary Across All Experiments}
\label{sec:summary}

\begin{table}[H]
\centering
\renewcommand{\arraystretch}{1.15}
\begin{tabular}{clll}
\toprule
\textbf{Phase} & \textbf{Approach} & \textbf{Best result} & \textbf{Status} \\
\midrule
BinFill & PPO + PPO+ICM, 250k steps & 0/20 SR (both) & done \\
\midrule
1 & V1 sweep (4 cells $\times$ 3 seeds $\times$ 1.5M) & 0/20 SR; max\_subgoal=0 & done \\
2 & BC pipeline (80 expert ep, loss 0.006) & 0/20 train AND test & done \\
3 & BC warm-start (V1/V2/V3) & 0/20 & done \\
4 & V4 + BC + aux=0.5 + low-std (5M target) & stuck at $-15$ floor & killed @850k \\
5 & DAgger v1 (4 iters, $K$=15) & 0/20 (polluted dataset) & done \\
6 & DAgger v2 (8 iters, $K$=5, fresh) & 0/20 through iter 5 & killed @iter 6 \\
7 & DAPG (5M $\times$ 3 seeds, all fixes) & stuck at $-15$ floor @590k & killed \\
\bottomrule
\end{tabular}
\caption{All experiments. Headline: \textbf{0\% eval SR} across every phase,
every baseline cell, every seed.}
\end{table}

\begin{finding}
\textbf{Bottom-line conclusion.}
Enhanced PPO with memory (LSTM) + curiosity (RND) + behavior cloning + DAgger
+ demo-augmented loss collapses to 0\% on PickXtimes within our 5M-step compute
budget ($\sim$165 Titan RTX-hours). This replicates the MIKASA-Robo (ICLR 2025)
finding that PPO/PPO-LSTM collapses on hard sparse-reward long-horizon manipulation
at standard scales.
\end{finding}

% ============================================================
\section{Engineering Fixes (Kept in Code)}
\label{sec:engineering}

\begin{table}[H]
\centering
\begin{tabular}{p{4.5cm}p{5cm}p{4cm}}
\toprule
\textbf{Issue} & \textbf{Fix} & \textbf{File} \\
\midrule
SAPIEN 3.0.3 + torch 2.9 GPU sim typestr error
& Force \texttt{ROBOMME\_SIM\_BACKEND=physx\_cpu} in all sbatch
& all \texttt{scripts/cluster/*.sbatch} \\
\addlinespace
Eval-time action clip (unbounded recurrent Gaussian $\to$ OOD env step)
& \texttt{np.clip(action, -1, 1)} before \texttt{env.step}
& \texttt{train/evaluate\_trained.py}, \texttt{scripts/eval\_*.py} \\
\addlinespace
Cluster slurm log dir wiped by rsync \texttt{-{}-delete}
& \texttt{-{}-exclude='logs/'} in sync script
& \texttt{scripts/cluster/sync\_to\_cluster.sh} \\
\addlinespace
BC \texttt{state\_dict} saved on GPU $\to$ load failure
& Save on CPU via \texttt{\{k: v.cpu() for k, v in ...\}}
& \texttt{scripts/bc\_pretrain\_pickxtimes.py} \\
\addlinespace
BC dataset OOM at 32G slurm allocation
& Stream per-episode to \texttt{.npz}; load slice on demand
& \texttt{scripts/collect\_bc\_pickxtimes.py} \\
\addlinespace
DAgger query captured ``stay-put + open-gripper'' no-ops
& Patch \texttt{planner.open/close\_gripper} to no-ops before \texttt{solve}
& \texttt{scripts/dagger\_iter.py:\allowbreak get\_first\_expert\_action} \\
\addlinespace
\texttt{eval\_subgoal} failed to load BC \texttt{.pt} state\_dicts
& Auto-detect \texttt{.pt} vs \texttt{.zip}; fresh model + \texttt{load\_state\_dict}
& \texttt{scripts/eval\_subgoal.py} \\
\bottomrule
\end{tabular}
\caption{Engineering fixes shipped and kept in the codebase.}
\end{table}

% ============================================================
\section{New Code Shipped}
\label{sec:code}

\begin{table}[H]
\centering
\begin{tabular}{ll}
\toprule
\textbf{Category} & \textbf{Files} \\
\midrule
Reward shaping (4 variants)
& \texttt{train/rewards/pickxtimes\{,\_v2,\_v3,\_v4\}.py} \\
RND module
& \texttt{train/models/rnd.py} \\
PPO+RND algo
& \texttt{train/algos/ppo\_with\_rnd.py} \\
RecurrentPPO+RND+BC-aux algo
& \texttt{train/algos/recurrent\_ppo\_with\_rnd.py} \\
Trainers (3 new)
& \texttt{train/train\_ppo\_rnd.py}, \texttt{train\_ppo\_recurrent.py},
  \texttt{train\_ppo\_recurrent\_rnd.py} \\
BC pipeline
& \texttt{scripts/collect\_bc\_pickxtimes.py},
  \texttt{scripts/bc\_pretrain\_pickxtimes.py} \\
DAgger pipeline
& \texttt{scripts/dagger\_iter.py} \\
Diagnostic evals
& \texttt{scripts/eval\_\{subgoal,stochastic,bc\_only,bc\_only\_train\}.py} \\
Aggregation
& \texttt{scripts/aggregate\_pickxtimes\_results.py} \\
Cluster sweep submitters
& \texttt{scripts/cluster/submit\_\{pickxtimes\_sweep,pickxtimes\_v2,}\\
& \quad\texttt{bc\_warmstart\_sweep,b6\_bc\_v2,b6\_bc\_v3,}\\
& \quad\texttt{b6\_bc\_v3\_lowent,5m\_bc\_v4,5m\_bc\_v4\_chain,dapg\_5m\}.sh} \\
Eval extensions
& \texttt{train/evaluate\_trained.py} (recurrent path + action-clip fix) \\
\bottomrule
\end{tabular}
\caption{New files added ($\approx$2{,}500 LOC total).}
\end{table}

% ============================================================
\section{Recommendations for Next Steps}
\label{sec:nextsteps}

In decreasing order of expected impact based on what the ablation has isolated:

\subsection{Strongest Option: Offline RL (IQL / CQL)}

\textbf{Why:} bypasses exploration entirely. A teammate's IQL result on this same
task is reported as non-zero, confirming the 42,821-transition BC dataset itself
is sufficient --- the bottleneck is the online RL paradigm, not data quality.

\begin{itemize}[nosep]
  \item The BC dataset is already collected and on disk at 5.9 GB
        (\texttt{scripts/collect\_bc\_pickxtimes.py} output).
  \item IQL (Kostrikov et al., 2021) or CQL (Kumar et al., 2020) train a
        $Q$-function entirely from offline demonstrations; policy extraction via
        advantage-weighted regression is more robust to compounding errors
        than BC's per-state MSE.
  \item \textbf{Implementation note}: the existing CNN extractor + policy class
        can be reused; only the training loop changes.
\end{itemize}

\subsection{Action Chunking (ACT-Style)}

Predicting $K=8$ actions at once reduces the effective horizon from 1500 to
$\approx$187 steps. This approach has been combined with PPO in recent manipulation
papers and addresses both F2 (exploration) and F3 (compounding errors) simultaneously.

\subsection{Hierarchical Action Space}

Expose planner primitives (reach, grasp, lift, drop, button-press) as the action
space; RL learns a meta-controller. Exploration drops from ``find a coordinated
8-D sequence over 1500 steps'' to ``pick the right one of 5 primitives at
$\sim$5 decision points.'' The motion planner (\texttt{PandaArmMotionPlanningSolver})
is already in the codebase and used for BC data collection.

\subsection{DAgger at Much Greater Scale}

DAgger v2 (8 iters) suggests BC retraining is the bottleneck — not the
expert-query mechanism itself. Running 50+ iterations with $K=2$ might
overcome it. Requires fixing the OOM issue (switch from pickle to HDF5 or
\texttt{zarr} for the aggregated dataset to avoid the 4-GiB pickle protocol 3
string limit).

\subsection{Reward V5 — Expert-Trajectory Distance}

Use the BC dataset as a reference set. Per-step reward $= -\min_{(s^*, a^*)}
\|s_t - s^*\|$ provides a dense gradient toward the expert manifold without
the proximity-hack failure mode (signal is ``be on expert trajectory,'' not
``be near object'').

\subsection{5–25M Training Steps}

MIKASA-Robo's positive results used 5–25M steps. Any of the above approaches
may still require this budget.

\subsection{Curriculum}

Train exclusively on \texttt{seed \% 3 == 0} (easy, $N=1$ cycle only) until
the first success, then add medium and hard. Reduces the sparse-reward
exploration problem by $\sim$3$\times$.

\subsection{What Is Explicitly Not Worth Trying}

\begin{itemize}[nosep]
  \item More ICM / RND / RIDE permutations on top of existing RecurrentPPO
        (already tested all combinations; memory and curiosity both fail).
  \item Additional reward variants V5–VN with any proximity term (all proximity
        formulations are exploitable at long horizons).
  \item BC warm-start alone without changing the RL loss structure (phases 3–7
        show it provides no eval SR improvement).
  \item Multi-task RL training (open research problem, out of scope).
\end{itemize}

% ============================================================
\section{Cluster Replication Notes}
\label{sec:replication}

\begin{enumerate}
  \item \textbf{Always export} \texttt{ROBOMME\_SIM\_BACKEND=physx\_cpu} before
        any training or eval script. GPU sim crashes on SAPIEN 3.0.2 + torch 2.9.

  \item \textbf{Action clipping}: eval scripts must call
        \texttt{np.clip(action, -1, 1)} before \texttt{env.step}.
        RecurrentPPO uses an unbounded Gaussian policy; without clipping, the
        deterministic-eval mean can step OOD joint targets.

  \item \textbf{rsync}: always include \texttt{-{}-exclude='logs/'} when syncing
        to cluster with \texttt{-{}-delete}. The \texttt{logs/} directory
        contains slurm output; deleting it causes exit-120 startup failures.

  \item \textbf{BC dataset streaming}: load the BC dataset per-episode from
        \texttt{.npz} files, not all-at-once. Full dataset is 5.9 GB; loading
        into RAM exceeds 32 GB slurm allocations.

  \item \textbf{Conda env}: \texttt{/vision/u/jevon/miniconda3/envs/robomme/}
        on the Stanford cluster. Contains all pinned versions in the table above.
\end{enumerate}

% ============================================================
\appendix

\section{File Pointers}
\label{app:files}

\begin{table}[H]
\centering
\begin{tabular}{ll}
\toprule
\textbf{Artifact} & \textbf{Path} \\
\midrule
BinFill eval report & \texttt{runs/eval\_results/REPORT.md} \\
PickXtimes findings & \texttt{runs/eval\_results/PICKXTIMES\_REPORT.md} \\
Full experiment log & \texttt{runs/eval\_results/EXPERIMENT\_LOG.md} \\
Next-steps plan & \texttt{runs/eval\_results/NEXT\_STEPS.md} \\
BinFill eval JSON & \texttt{runs/eval\_results/ppo\_v4\_BinFill\_official.json} \\
Training logs & \texttt{runs/<run>/train.log} \\
BC dataset & \texttt{data/bc\_pickxtimes/} (on cluster; 5.9 GB) \\
DAgger v2 dataset & \texttt{data/dagger\_v2/} (on cluster; iters 1–5) \\
Cluster conda env & \texttt{/vision/u/jevon/miniconda3/envs/robomme/} \\
\bottomrule
\end{tabular}
\end{table}

\section{Reward Function Code Sketches}

\paragraph{V1 (proximity-exploitable).}

\begin{lstlisting}[language=python]
def compute_reward(obs, info):
    d = tcp_to_cube_distance(obs)
    r_prox  = 0.05 / (1 + d)       # exploitable: agent hovers near cube
    r_grasp = 0.05 * is_grasping(info)
    r_sub   = 2.0 * subgoal_advance(info)
    r_pick  = 1.0 * pickup_detected(info)
    r_term  = +10 if info["success"] else (-5 if info["fail"] else 0)
    return r_prox + r_grasp + r_sub + r_pick + r_term - 0.01
\end{lstlisting}

\paragraph{V4 (no proximity, pure task signal).}

\begin{lstlisting}[language=python]
def compute_reward(obs, info):
    # alpha = 0: no proximity term
    r_grasp = 0.5 * is_grasping(info)
    r_sub   = 10.0 * subgoal_advance(info)
    r_pick  = 5.0  * pickup_detected(info)
    r_term  = +100 if info["success"] else (-10 if info["fail"] else 0)
    return r_grasp + r_sub + r_pick + r_term - 0.01
\end{lstlisting}

\end{document}
